\section{Generador de carga y rendimiento}

El generador de carga vive en \codigo{src/main/java/com/escom/banco/carga/GeneradorCarga.java}.
Es una herramienta de linea de comandos en Java puro, sin dependencias nuevas,
que golpea al nodo lider durante un minuto con una mezcla de 80\% lecturas y
20\% transferencias concurrentes, y al terminar comprueba que la suma total de
saldos no haya cambiado (invariante de conservacion del dinero).

\subsection{Modelo de carga}

La herramienta primero se autentica como el usuario \codigo{carga}, lee el numero
de cuentas y el saldo total inicial desde \codigo{/stats}, y exige al menos dos
cuentas para poder transferir. Cada peticion se decide al azar: con probabilidad
del 80\% es una lectura \codigo{GET /api/accounts/\{id\}} y con el 20\% restante
es una transferencia \codigo{POST /api/transactions/transfer} entre dos cuentas
elegidas al azar. Los montos se generan en centavos y se serializan con
\codigo{BigDecimal} para no perder exactitud al pasarlos a JSON.

\begin{lstlisting}[language=Java,caption={Decision 80/20 por peticion (lanzarUna)}]
boolean lectura = rnd.nextInt(100) < 80;
if (lectura) {
    int id = 1 + rnd.nextInt(cuentas);
    req = HttpRequest.newBuilder(URI.create(base + "/api/accounts/" + id))
            .header("Authorization", autorizacion).GET().build();
} else {
    int origen = 1 + rnd.nextInt(cuentas);
    int destino = 1 + rnd.nextInt(cuentas);
    long montoCent = 1 + rnd.nextLong(cfg.montoMaxCent());
    // ...
}
\end{lstlisting}

La parametrizacion de una corrida se concentra en el record \codigo{Config}, que
distingue las peticiones en vuelo del numero de clientes HTTP independientes.

\begin{lstlisting}[language=Java,caption={Parametros de una corrida (Config)}]
public record Config(String host, int puerto, int segundos, int enVuelo,
                     long montoMaxCent, int clientes) {
    /** Conveniencia: un solo HttpClient (comportamiento original). */
    public Config(String host, int puerto, int segundos, int enVuelo, long montoMaxCent) {
        this(host, puerto, segundos, enVuelo, montoMaxCent, 1);
    }
}
\end{lstlisting}

\subsection{Motor asincrono con ventana en vuelo}

En lugar de dedicar un hilo del sistema operativo por peticion (un envio
bloqueante a la vez), el generador mantiene una ventana de N peticiones EN VUELO
usando \codigo{HttpClient.sendAsync()} y un \codigo{Semaphore} que limita cuantas
peticiones pueden estar pendientes a la vez. Cada permiso se adquiere antes de
lanzar la peticion y se libera en el callback \codigo{whenComplete}, de modo que
una sola maquina sostiene mucha mas concurrencia sin gastar N hilos.

\begin{lstlisting}[language=Java,caption={Bucle del motor con semaforo (motor)}]
private void motor(HttpClient cliente, int cuentas, int ventana, long fin) {
    Semaphore permisos = new Semaphore(ventana);
    try {
        while (System.nanoTime() < fin) {
            permisos.acquire();
            if (System.nanoTime() >= fin) { permisos.release(); break; }
            lanzarUna(cliente, cuentas, permisos);
        }
        permisos.acquire(ventana); // drenar las peticiones en vuelo de este motor
    } catch (InterruptedException e) {
        Thread.currentThread().interrupt();
    }
}
\end{lstlisting}

Al terminar la fase de envio, el motor adquiere la ventana completa
(\codigo{permisos.acquire(ventana)}) para drenar las peticiones que aun siguen en
vuelo antes de devolver el control. El conteo de exitos y rechazos se acumula con
\codigo{LongAdder}, una estructura concurrente que escala mejor que un contador
unico cuando muchos callbacks incrementan a la vez.

\subsection{Parametro clientes (K) y verificacion del invariante}

El parametro \codigo{clientes} (K) crea K instancias de \codigo{HttpClient}
independientes, cada una con su propio hilo selector de I/O. La ventana total de
peticiones en vuelo se reparte entre ellas (\codigo{ventana = enVuelo / k}), de
modo que UNA sola maquina generadora no se tope con el limite de un unico
selector. El \codigo{HttpClient} del campo compartido se usa solo para
autenticarse y leer \codigo{/stats}; la carga real la llevan los K clientes
creados por hilo.

\begin{lstlisting}[language=Java,caption={Reparto de la ventana entre K clientes (ejecutar)}]
int k = Math.max(1, cfg.clientes());
int ventana = Math.max(1, cfg.enVuelo() / k); // ventana POR cliente
long fin = System.nanoTime() + cfg.segundos() * 1_000_000_000L;

Thread[] hilos = new Thread[k];
for (int i = 0; i < k; i++) {
    HttpClient cliente = HttpClient.newBuilder()
            .version(HttpClient.Version.HTTP_1_1)
            .connectTimeout(Duration.ofSeconds(5)).build();
    hilos[i] = new Thread(() -> motor(cliente, cuentas, ventana, fin), "motor-" + i);
}
for (Thread h : hilos) h.start();
for (Thread h : hilos) h.join();
\end{lstlisting}

Una vez que todos los hilos terminan, se vuelve a leer el saldo total y se
compara con el inicial. Si coinciden, el dinero se conservo exacto y se imprimen
las transacciones y lecturas exitosas; si difieren, se reporta la inconsistencia
indicando el delta en centavos y se sale con codigo de error.

\begin{lstlisting}[language=Java,caption={Verificacion del invariante (main)}]
if (r.invarianteOk()) {
    System.out.println("INVARIANTE OK: la suma total se conservo exacta.");
} else {
    System.out.printf("INCONSISTENCIA: el total cambio en %s (centavos: %d).%n",
            centavosAJson(r.deltaCent()), r.deltaCent());
    System.exit(1);
}
\end{lstlisting}

La interfaz de linea de comandos es
\codigo{host puerto [segundos] [enVuelo] [clientes]}, con valores por omision
de 60 segundos, 512 peticiones en vuelo y 1 cliente:

\begin{lstlisting}[language=none,caption={Uso por linea de comandos}]
Uso: GeneradorCarga <host> <puerto> [segundos=60] [enVuelo=512] [clientes=1]
\end{lstlisting}

\subsection{Resultados}

La progresion de rendimiento medida (banco fijo en una \codigo{e2-standard-4})
muestra como se fue moviendo el cuello de botella. La primera version sobre
\codigo{com.sun.net.httpserver} topaba en unas 16,000 ops/s por el dispatcher de
un solo hilo. Al pasar a un servidor NIO y generar carga desde una
\codigo{e2-standard-8} se alcanzaron 25,000 a 28,000 ops/s; con NIO y un generador
\codigo{c3-standard-22} usando 1 \codigo{HttpClient} se llego a unas 60,000 ops/s;
y con \codigo{c3-standard-22} y K igual a 2 o 3 \codigo{HttpClient} se llego a
unas 110,000 ops/s, saturando la CPU del banco al 96 a 98\%.

Los numeros oficiales por escenario se midieron en 60 s, con un generador
\codigo{c3-standard-22} y K igual a 3, en la configuracion de concurso con Cloud
Storage y replicacion activos. El invariante de conservacion se mantuvo exacto en
los tres escenarios.

\begin{tablaglab}{L{4.5cm} R{2.6cm} R{2.8cm} R{2.5cm}}
\thead{Escenario} & \thead{Lecturas} & \thead{Transferencias} & \thead{Ops/s aprox.} \\ \midrule
Nodos 1+2+3 & 5,523,490 & 1,379,928 & 115,057 \\
Nodos 1+2   & 6,231,510 & 1,558,375 & 129,831 \\
Nodo 1      & 6,353,653 & 1,588,544 & 132,370 \\
\end{tablaglab}

El sobrecosto (overhead) de replicacion es de aproximadamente 13\% con dos
replicas y de aproximadamente 2\% con una sola replica. El puntaje del concurso
pondera las transferencias 4x sobre las lecturas y suma los tres escenarios.
